\documentclass[11pt]{article}
\usepackage[T1]{fontenc}
\usepackage[utf8]{inputenc}
\usepackage{lmodern}
\usepackage{microtype}
\usepackage{amsmath,amssymb}
\usepackage{array}
\usepackage{geometry}
\usepackage{tikz}
\usetikzlibrary{calc,intersections,arrows.meta,decorations.pathreplacing}
\geometry{paperwidth=6.6in,paperheight=9.6in,inner=0.68in,outer=0.68in,top=0.72in,bottom=0.78in,marginparwidth=0.55in,marginparsep=0.06in}
\setlength{\parindent}{1.2em}
\setlength{\parskip}{0.18em}
\setcounter{secnumdepth}{0}
\newcommand{\paperhead}[3]{%
  \clearpage
  \begin{center}
  {\Large\bfseries #1.}\par\vspace{0.85em}
  {\large\scshape #2}\par\vspace{0.6em}
  {\itshape [#3]}\par
  \end{center}\vspace{0.4em}}
\newcommand{\subpaperhead}[1]{\begin{center}{\large\scshape #1}\end{center}}
\newcommand{\proposition}[1]{\vspace{0.9em}\begin{center}{\scshape Proposition #1.}\end{center}\vspace{0.2em}}
\newcommand{\sourcepage}[1]{\marginpar{\scriptsize p.~#1}}
\newcommand{\tagtext}[1]{\tag{#1}}
\newcommand{\dd}{\,d}
\newcommand{\braced}[1]{\left\{#1\right\}}
\newcommand{\cosec}{\operatorname{cosec}}
\newcommand{\figcaption}[1]{\par\smallskip\centerline{\small #1}\smallskip}

\newcommand{\theoremhead}[1]{\vspace{0.75em}\begin{center}{\scshape Theorem #1.}\end{center}\vspace{0.15em}}
\newcommand{\corhead}[1]{\textsc{Cor. #1.}}
\newcommand{\art}[1]{\par\smallskip\noindent Art.~(#1).\quad}
\newcommand{\SPD}{\mathit{SPD}}
\newcommand{\QDP}{\mathit{QDP}}
\newcommand{\PD}{\mathit{PD}}
\newcommand{\SR}{\mathit{SR}}
\newcommand{\rowpair}[2]{\left\{\begin{matrix}#1\\#2\end{matrix}\right\}}
\newcommand{\rowparen}[2]{\left(\begin{matrix}#1\\#2\end{matrix}\right)}
\newcommand{\arb}{\mathfrak{S}}
\newcommand{\zprod}{\zeta}
\newcommand{\pgram}[1]{\overline{#1}}
\newcommand{\dotter}[1]{\dot{#1}}
\newcommand{\tworow}[2]{\left(\begin{matrix}#1\\#2\end{matrix}\right)}
\newcommand{\eps}{\varepsilon}

\newcommand{\bd}[2]{\left\{\begin{matrix}#1\\#2\end{matrix}\right\}}
\newcommand{\bds}[2]{\left\{\begin{smallmatrix}#1\\#2\end{smallmatrix}\right\}}
\newcommand{\matdet}[1]{\left|\begin{matrix}#1\end{matrix}\right|}
\begin{document}

\begin{center}
{\Large James Joseph Sylvester}\par\vspace{0.25em}
{\large Collected Mathematical Papers, Vol. I}\par\vspace{0.25em}
{\normalsize New work: Paper 43, book pp.~328--347}\par
\end{center}
\thispagestyle{empty}
\vspace{0.75em}


\paperhead{43}{On the Principles of the Calculus of Forms}{Cambridge and Dublin Mathematical Journal, VII. (1852), pp.~179--217}

\subpaperhead{Part I. \quad Section IV. \quad Reciprocity, also Properties and Analogies of certain Invariants, \&c.}

\sourcepage{328}
It will hereafter be found extremely convenient to represent all systems of variables cogredient with the original system in the primitive form by letters of the Roman, and all contragredient systems by letters of the Greek alphabet; the rules for concomitance may then be applied without paying any regard to the distinction between the direction of the march of the substitutions, the variables at the close of each operation as it were telling their own tale in respect of being cogredients or contragredients. This distinction has not (as it should have) been uniformly observed in the preceding sections; as, for instance, in the notation for emanants which have been derived by the application of the symbol
\[
\left(\xi\frac{d}{dx}+\eta\frac{d}{dy}+\&c.\right)^2,
\]
instead of the more appropriate one
\[
\left(x'\frac{d}{dx}+y'\frac{d}{dy}+\&c.\right)^2.
\]

The observations in this section will refer exclusively to points of doctrine which have been started in the preceding sections in such order as they more readily happen to present themselves. And, first, as to some important applications of the reciprocity method referred to in Notes (6) and (8) of the Appendix [pp.~325, 327 above].

The practical application of this method will be found greatly facilitated by the rule that $x,y,z$, \&c. may always in any combination of concomitants be replaced respectively by
\[
\frac{d}{d\xi},\quad \frac{d}{d\eta},\quad \frac{d}{d\zeta},\quad \&c.,
\]
and \emph{vice versâ}. I shall apply this prolific principle of reciprocity to elucidate some of the properties and relations of Aronhold's $S$ and $T$, and certain other kindred forms. This $S$ and $T$ are the quartinvariant and sextinvariant respectively of a cubic of three variables. I give the names of $s$ and $t$ to the quadrinvariant and cubinvariant of the quartic function of two variables. Furthermore, whoever will consider attentively the remarks made in Section II. of the foregoing relative to reciprocal polars, will apprehend without any difficulty that to every invariant of a function of any degree of any number of variables will

\sourcepage{329}
correspond a contravariant of a function of the same degree of variables one more in number, and that between such invariants, whatever relations exist expressed independently of all other quantities, precisely the same relations must exist between the corresponding contravariants. Thus, then, to $s$ and $t$ the two invariants of $(x,y)^4$ will correspond two contravariants $\sigma$ and $\tau$ of $(x,y,z)^4$, and to $S$ and $T$ the two invariants of $(x,y,z)^3$ will correspond $\Sigma$ and $\mathfrak{S}$ two contravariants of $(x,y,z,t)^3$. Calling $r$ the resultant of $(x,y)^4$, $R$ the resultant of $(x,y,z)^3$, $\rho$ the polar reciprocal, or, more briefly, the reciprocant of $(x,y,z)^4$, and $(R)$ the reciprocant of $(x,y,z,t)^3$, we have the following equations (presuming that all the quantities are previously affected with the proper numerical multipliers), namely
\[
r=s^3+t^2,\qquad \rho=\sigma^3+\tau^2,
\]
\[
R=S^3+T^2,\qquad (R)=\Sigma^3+\mathfrak{S}^2.
\]

I propose in this First Annotation to point out the remarkable analogies which exist between the modes of generating the four pairs of quantities $s,t$, \&c., the functions severally corresponding to which I shall call $u,\omega,U,\Omega$. The Hessian corresponding to any of these functions will be denoted by an $H$ prefixed, and when we have to consider, not the pure Hessian, but the matrix formed from it by adding a vertical and horizontal border of variables, the same in number but contragredient to the variable of the function (as, for instance, the Hessian of $u$ bordered with $\xi,\eta$ horizontally and vertically, or of $U$ with $\xi,\eta,\zeta$), then I shall denote the result by the ruled symbol $\overline H$, and if there be occasion to add two borders, as $\xi,\eta,\zeta;\xi',\eta',\zeta'$, both repeated in the horizontal and vertical directions, the result will be typified by the doubly ruled $H$, $\overline{\overline H}$.

Now, in the first place, as observed by me in Note (8) of the Appendix in the last number; if we call the coefficients of $U$ (10 in number) $a,b,c,d$, \&c., we have
\[
S=\overline H\!\left(
\frac{d}{d\xi},\frac{d}{d\eta},\frac{d}{d\zeta};
\frac{d}{dx},\frac{d}{dy},\frac{d}{dz}\right)
\overline H(x,y,z;\xi,\eta,\zeta),
\]
also
\[
T=\frac{dS}{da}\frac{d^3H}{dx^3}
+\frac{dS}{db}\frac{d^3H}{dx^2dy}
+\frac{dS}{dc}\frac{d^3H}{dx^2dz}+\&c.
\]
I will now add the further important relation
\[
S^2=\frac{dT}{da}\frac{d^3H}{dx^3}
+\frac{dT}{db}\frac{d^3H}{dx^2dy}
+\frac{dT}{dc}\frac{d^3H}{dx^2dz}+\&c.\footnote{It will be found hereafter convenient to designate contravariants formed in this manner from invariants as Evects of such invariants or contravariants, and according to the number of times that such process of derivation is applied, 1st, 2nd, 3rd, \&c. evects. Such evects form a peculiar class, and when considered generally, without reference to the base to which they refer, they may be termed evectants. Evectants will be again distinguishable according as their base is an invariant simply or a contravariant. Perhaps the terms pure and affected evectants may serve to mark this distinction.}
\]

\sourcepage{330}
so that it will be observed if all the derivatives of $S$ are zero, $T$ is zero, and \emph{vice versâ}.

Precisely in the same way, using $h$ and $\overline h$ to denote respectively the Hessian of $u$ and the same bordered with $\xi,\eta$, we have
\[
s=\overline h\!\left(\frac{d}{d\xi},\frac{d}{d\eta};
\frac{d}{dx},\frac{d}{dy}\right)\overline h(x,y;\xi,\eta),
\]
\[
t=\frac{ds}{da}\frac{d^4h}{dx^4}
+\frac{ds}{db}\frac{d^4h}{dx^3dy}
+\frac{ds}{dc}\frac{d^4h}{dx^2dy^2}+\&c.,
\]
\[
s^2=\frac{dt}{da}\frac{d^4h}{dx^4}
+\frac{dt}{db}\frac{d^4h}{dx^3dy}
+\frac{dt}{dc}\frac{d^4h}{dx^2dy^2}+\&c.
\]

Again, taking $(\overline{\overline H})$ the second bordered Hessian of $\Omega$; that is, $\Omega$ bordered as well horizontally as vertically with the double lines and columns $\xi,\eta,\zeta,\theta;\xi',\eta',\zeta',\theta'$,
\[
\begin{aligned}
\Sigma&=(\overline{\overline H})\!\left(
\frac{d}{d\xi},\frac{d}{d\eta},\frac{d}{d\zeta},\frac{d}{d\theta};
\frac{d}{dx},\frac{d}{dy},\frac{d}{dz},\frac{d}{dt};
\xi',\eta',\zeta',\theta'\right)\\
&\qquad{}\times
(\overline{\overline H})(x,y,z,t;\xi,\eta,\zeta,\theta;\xi',\eta',\zeta',\theta'),
\end{aligned}
\]
\[
\mathfrak{S}=\frac{d\Sigma}{da}\frac{d^3\overline H}{dx^3}
+\frac{d\Sigma}{db}\frac{d^3\overline H}{dx^2dy}
+\frac{d\Sigma}{dc}\frac{d^3\overline H}{dx^2dz}
+\frac{d\Sigma}{dd}\frac{d^3\overline H}{dx^2dt}+\&c.,
\]
\[
\Sigma^2=\frac{d\mathfrak{S}}{da}\frac{d^3\overline H}{dx^3}
+\frac{d\mathfrak{S}}{db}\frac{d^3\overline H}{dx^2dy}+\&c.
\]
In like manner again
\[
\begin{aligned}
\sigma&=(\overline h)\!\left(
\frac{d}{d\xi},\frac{d}{d\eta},\frac{d}{d\zeta};
\frac{d}{dx},\frac{d}{dy},\frac{d}{dz};
\xi',\eta',\zeta'\right)\\
&\qquad{}\times
\overline h\{x,y,z;\xi,\eta,\zeta;\xi',\eta',\zeta'\},
\end{aligned}
\]
\[
\tau=\frac{d\sigma}{da}\frac{d^4(\overline h)}{dx^4}+\&c.,
\qquad
\sigma^2=\frac{d\tau}{da}\frac{d^4\overline h}{dx^4}+\&c.
\]

$\sigma$ and $\tau$ are the same quantities as are calculated by Mr Salmon, in his inestimable work \emph{On Higher Plane Curves}, but are there expressed under the names of $S$ and $T$, with the sole difference that in place of $x,y,z$, used by Mr Salmon, the contragredient variables $\xi,\eta,\zeta$ are used in the expressions above. Mr Salmon has also pointed out to me that $\sigma$ may be obtained by operating with
\[
\left(\xi^4\frac{d}{da}+\xi^3\eta\frac{d}{db}
+\xi^2\zeta^2\frac{d}{dc}+\&c.\right)
\]
directly upon $I$ a cubic invariant of the function $u$, or $(x,y,z)^4$. This $I$ is no other than the simple commutant obtained by operating upon $u$ with the commutantive symbol formed by taking four times over the line
\[
\frac{d}{dx},\quad \frac{d}{dy},\quad \frac{d}{dz},
\]
agreeable to the remark made in the third section that

\sourcepage{331}
every function of an even degree of $n$ variables possesses an invariant of the $n$th order in extension of Mr Cayley's observation that every such function of two variables possesses a quadrinvariant, that is an invariant of the second order.

I need hardly remark that $\sigma$ is of 2 dimensions in the coefficients and of 4 in the contragredient variables, $\tau$ of 3 in the coefficients and of 5 in the contragredients, $\Sigma$ of 4 in the constants and 4 in the contragredients, $\mathfrak{S}$ of 6 in the constants and 6 in the contragredients, or that the single-bordered Hessians of $u$ and $U$ and the double-bordered Hessians of $\omega$ and $\Omega$ are each of them quadratic in respect of the $x$ \&c. as well as of the $\xi$ \&c. systems.

If the right numerical factors be attributed to $S,T$, Aronhold has shown that
\[
H\{H(U)\}+T\cdot H(U)+S^2U=0,
\]
and in my paper in the last May Number\footnote{p.~192 above.}, I gave the equation
\[
h\{h(u)\}+s\cdot h(u)+tu=0.
\]
I think it highly probable that it will be found that the analogous equations obtain, namely
\[
\overline H\{\overline H(\Omega)\}
+\mathfrak{S}\cdot\overline H(\Omega)+\Sigma^2\Omega=0,
\]
\[
\overline h\{\overline h(\omega)\}+\sigma\cdot\overline h(\omega)+\tau\omega=0.
\]
These remarkable equations, if verified (of which I can scarcely doubt), will be most powerful aids to the dissection of the forms $\omega,\Omega$, and thereby to the detection of the fundamental properties of curves of the fourth and surfaces of the third degree, of which at present so little is known. It will have been observed that in the preceding developments the contravariants of $\omega$ and $\Omega$ were derived in precisely the same way from $\omega$ and $\Omega$ as the corresponding invariants of $u$ and $U$ from $u$ and $U$, with the sole difference that the Hessian used in the two latter cases is replaced by a single-bordered Hessian in the two former cases, and a single-bordered Hessian in the two latter by a double-bordered Hessian in the two former. The analogies are not even yet stated exhaustively; for it will be remembered (as shown in the third section), that $T$ and $S$ can be derived directly and concurrently by means of operating with the commutantive symbol
\[
\left.
\begin{matrix}
\dfrac{d}{dx},&\dfrac{d}{dy},&\dfrac{d}{dz}\\
\dfrac{d}{dx},&\dfrac{d}{dy},&\dfrac{d}{dz}\\
\dfrac{d}{d\xi},&\dfrac{d}{d\eta},&\dfrac{d}{d\zeta}\\
\dfrac{d}{d\xi},&\dfrac{d}{d\eta},&\dfrac{d}{d\zeta}
\end{matrix}
\right\}
\quad\text{upon }\overline H(U)+\lambda(x\xi+y\eta+z\zeta)^2,
\]

\sourcepage{332}
which gives a result of the form $m(\lambda^3+S\lambda+T)$, $m$ being a number; and I conjecture that if
\[
\left.
\begin{matrix}
\dfrac{d}{dx},&\dfrac{d}{dy},&\dfrac{d}{dz},&\dfrac{d}{dt}\\
\dfrac{d}{dx},&\dfrac{d}{dy},&\dfrac{d}{dz},&\dfrac{d}{dt}\\
\dfrac{d}{d\xi},&\dfrac{d}{d\eta},&\dfrac{d}{d\zeta},&\dfrac{d}{d\theta}\\
\dfrac{d}{d\xi},&\dfrac{d}{d\eta},&\dfrac{d}{d\zeta},&\dfrac{d}{d\theta}
\end{matrix}
\right\}
\]
be made to operate upon
\[
\overline{\overline H}\Omega+\lambda(x\xi+y\eta+z\zeta+t\theta)^2,
\]
and the result be put under the form
\[
m(\lambda^4+A\lambda^3+B\lambda^2+C\lambda+D),
\]
that $A$ will be zero, $B$ and $C$ will be respectively $\Sigma$ and $\mathfrak{S}$, and perhaps $D$ (a contravariant, if it effectively exist, of 8 dimensions in the coefficients of $\Omega$, and of a like number in the contragredients $\xi',\eta',\zeta',\theta'$), also zero. But of the evanescence of $D$ I do not speak with any degree of assurance.

Mr Salmon has made an excellent observation to the effect that if we call $(\sigma)$ what $\sigma$ becomes when $\xi',\eta',\zeta'$ are replaced by
\[
\frac{d}{dx},\quad \frac{d}{dy},\quad \frac{d}{dz},
\]
$(\sigma)h(\omega)$ will represent a covariant to $\omega$ of $3+2$, that is, 5 dimensions in the coefficients, and of $6-4$, that is, of 2 dimensions in $x,y,z$, $h(\omega)$ being of 3 and 6 dimensions in these respectively, and $\sigma$ of 2 and 4 dimensions respectively in the same. Now these resulting dimensions 5 and 2 precisely agree with the form especially noticed by me in Note\footnote{p.~323 above.} (2) of the Appendix, where it was derived as one of a group by the method of unravelment. There can be little doubt that these two conics each of them indissolubly connected with every curve of the fourth degree are identical. The form $(\sigma)h(\omega)$ enables us to prove readily (thanks to Mr Salmon's calculation of $\sigma$, given in his \emph{Higher Plane Curves}, under the name of $S$) that this is a \emph{bonâ fide} existent conic.

For if we take a particular case of $\omega$, say
\[
\omega=a_1x^4+b_2y^4+c_3z^4+6dy^2z^2,
\]
we find
\[
h(\omega)=
\left|\begin{array}{ccc}
a_1x^2&0&0\\
0&b_2y^2+dz^2&dyz\\
0&dyz&c_3z^2+dy^2
\end{array}\right|
\]
\[
=a_1(b_2c_3+d^2)x^2y^2z^2+a_1b_2dx^2y^4+a_1c_3dx^2z^4,
\]

\sourcepage{333}
and $\sigma$ becomes
\[
a_1d\eta^2\zeta^2,
\]
and consequently $(\sigma)$ is
\[
a_1d\left(\frac{d}{dy}\right)^2\left(\frac{d}{dz}\right)^2;
\]
and therefore
\[
(\sigma)h(\omega)=4a_1^2d(b_2c_3+d^2)x^2,
\]
the conic here reducing to a pair of coincident straight lines. This example demonstrates that the conic is in general actually existent.

As I have said so much upon $S$ and $T$ it may not be irrelevant to state in this place how I obtained the conditions for $U$, the characteristic of the curve of the third degree becoming the characteristic of a conic and a straight line, that is breaking up into a linear and a quadratic factor, which Mr Salmon has inserted in the notes to his work above referred to. When $U$ is of this form it may obviously by linear transformations be expressed by $ax^3+6dxyz$, but when starting with the general form,
\[
a_1x^3+b_2y^3+c_3z^3+\&c.+6Dxyz,
\]
we form two contravariants from $S$ and $T$, to wit
\[
\left(\xi^3\frac{d}{da_1}+\eta^3\frac{d}{db_2}
+\zeta^3\frac{d}{dc_3}+\&c.+\xi\eta\zeta\frac{d}{dD}\right)S,\quad\text{say }S',
\]
\[
\left(\xi^3\frac{d}{da_1}+\eta^3\frac{d}{db_2}
+\zeta^3\frac{d}{dc_3}+\&c.+\xi\eta\zeta\frac{d}{dD}\right)T,\quad\text{say }T',
\]
and then make $a_1=a$, $D=d$, and all the other coefficients zero, it will easily be seen on examining the forms of $S$ and $T$, given by Mr Salmon, that $(S')$ and $(T')$ (the evectants of $S$ and $T$) become respectively
\[
4d^2\xi\eta\zeta,\qquad 31d^5\xi\eta\zeta;
\]
we have therefore $(T')+\lambda(S')=0$: and $(T')$ and $(S')$, although contravariantive to their primitive $U$, are covariantive with one another, so that $(T')+\lambda(S')=0$ is a persistent relation unaffected by linear transformations; it follows therefore that when $U$ is of, or reducible to, the form supposed,
\[
\frac{dS}{da_1}:\frac{dS}{db_2}:\frac{dS}{dc_3}:\&c.:\frac{dS}{dD}
=
\frac{dT}{da_1}:\frac{dT}{db_2}:\frac{dT}{dc_3}:\&c.:\frac{dT}{dD},
\]
which is the criterion given in the note referred to.\footnote{Mr Salmon has remarked to me that the two evectants $(S)$ and $(T)$ intersect in the nine cuspidal points of the polar reciprocal to the curve.}

I am also able to obtain these equations more directly by another method founded upon a New View of the Theory of Elimination, an account of which,

\sourcepage{334}
however, I must reserve for another occasion, but which, I may mention, serves to fix not merely the conditions, as in the ordinary restricted theory, that a given set of equations may be simultaneously satisfiable by some one system of values of the variables, but the \emph{conditions} that such set of equations may be simultaneously satisfiable by any given number of distinct systems of variables.

Mr Salmon has remarked to me to the effect that if in $\tau$ we write
\[
\frac{d}{dx},\quad \frac{d}{dy},\quad \frac{d}{dz},
\]
in place of the contragredients, and call $\tau$ so altered $(\tau)$, then $(\tau)h(\omega)$ will be an invariant of 6 dimensions in the coefficients of $\omega$. This sextinvariant I have little doubt is identical with that obtained by operating upon $\omega$ with the commutantive symbol
\[
\left.
\begin{matrix}
\left(\dfrac{d}{dx}\right)^2,&\dfrac{d}{dx}\dfrac{d}{dy},&
\left(\dfrac{d}{dy}\right)^2,&\dfrac{d}{dy}\dfrac{d}{dz},&
\left(\dfrac{d}{dz}\right)^2,&\dfrac{d}{dz}\dfrac{d}{dx}\\
\left(\dfrac{d}{dx}\right)^2,&\dfrac{d}{dx}\dfrac{d}{dy},&
\left(\dfrac{d}{dy}\right)^2,&\dfrac{d}{dy}\dfrac{d}{dz},&
\left(\dfrac{d}{dz}\right)^2,&\dfrac{d}{dz}\dfrac{d}{dx}
\end{matrix}
\right\}.
\]
This, like every other commutant of 2 lines only, is of course capable of being expressed under the form of an ordinary determinant, and the remark is not without interest, as showing how the proposition, known with respect to quadratic functions of any number of variables, namely of every such having an invariantive determinant, lends itself to the general case of functions of any even degree of any number of variables which also have always an invariantive determinant attached to them, of which the terms are simple coefficients of such functions. The only peculiarity (if it be one) of quadratic functions in this respect being that they have each but one invariant of such form and no other. In the case before us, if we write
\[
\begin{aligned}
\omega={}&a_1x^4+b_2y^4+c_3z^4+4a_2x^3y+4a_3x^3z+4b_1y^3x+4b_3y^3z\\
&+4c_1z^3x+4c_2z^3y+6dy^2z^2+6ez^2x^2+6fx^2y^2\\
&+12lx^2yz+12mxy^2z+12nxyz^2,
\end{aligned}
\]
the sextinvariant in question becomes representable under the form of the determinant
\[
\left|\begin{array}{cccccc}
a_1&a_2&f&l&e&a_3\\
a_2&f&b_1&m&n&l\\
f&b_1&b_2&b_3&d&m\\
l&m&b_3&d&c_2&n\\
e&n&d&c_2&c_3&c_1\\
a_3&l&m&n&c_1&e
\end{array}\right|.\footnote{This determinant is identical with the determinant formed by taking the second differential coefficients of the function and arranging in the usual manner the coefficients of the several powers and combinations of powers of the variables treated as if they were independent quantities.}
\]

\sourcepage{335}
Before quitting the subject of $S$ and $T$ the two invariants of the cubic function of 3 variables, or, as it may be termed, of the cubic curve, it may not be amiss to give the complete table which I have formed corresponding to all the singular cases which can befall such curve, which will be seen below to be eight in number; it is of the highest importance to push forward the advanced posts of geometry, and for this purpose to obtain the same kind of absolute power and authority over, and clear and absolute knowledge of, the properties and affections of cubic forms as have been already attained for forms of the second degree.

Let
\[
U=ax^3+4bx^2y+4cxz^2+\&c.
\]

(1) When $U$ has one double point
\[
S^3+T^2=0.
\]

(2) When $U$ has two double points, that is becomes a conic and right line
\[
\frac{dS}{da}\frac{dT}{db}-\frac{dS}{db}\frac{dT}{da}=0,\quad \&c.\ \&c.
\]

(3) When $U$ has a cusp
\[
S=0,\qquad T=0.
\]

(4) When $U$ has two coincident double points, that is, is a conic and a tangent line thereto, which comprises the two preceding cases in one,
\[
\frac{dT}{da}=0,\qquad \frac{dT}{db}=0,\quad \&c.,
\]
and also therefore
\[
S=0.
\]

(5) When $U$ becomes three right lines forming a triangle
\[
\frac{d^2S}{da\,db}\frac{d^2T}{dc\,de}
-\frac{d^2T}{da\,db}\frac{d^2S}{dc\,de}=0,\quad \&c.,
\]
where $a,b,c,e$ each represent any of the coefficients arbitrarily chosen, whether distinct or identical.

Another, and lower in degree system of equations, may be substituted for the above, obtained by affirming the equality of the ratios between the coefficients of $U$ and the corresponding coefficients of its Hessian.

(6) When $U$ represents a pencil of three rays meeting in a point
\[
\frac{dS}{da}=0,\qquad \frac{dS}{db}=0,\quad \&c.
\]
and also therefore
\[
T=0.
\]
Also in place of this system may be substituted the system obtained by taking all the coefficients of the Hessian zero.

\sourcepage{336}
(7) When $U$ becomes a line, and two other coincident lines,
\[
\frac{dS}{da}=0,\qquad \frac{dS}{db}=0,\quad \&c.
\]
and also
\[
\frac{d^2T}{da^2}=0,\qquad \frac{d^2T}{da\,db}=0,\quad \&c.
\]
I have not ascertained whether this second system necessarily implies the first; I rather think that it does not. In the preceding case also it would be interesting to show the direct algebraical connexion between the system formed by the coefficients of the Hessian and the system consisting of the first derivatives of $S$.

(8) When $U$ becomes a perfect cube representing three coincident right lines
\[
\frac{d^2S}{da^2}=0,\qquad \frac{d^2S}{da\,db}=0,\quad \&c.
\]
and
\[
\frac{d^2T}{da^2}=0,\qquad \frac{d^2T}{da\,db}=0,\quad \&c.
\]
The first of these systems of equations necessarily implies the equations
\[
\frac{dT}{da}=0,\quad \frac{dT}{db}=0,\quad \&c.,
\]
as is obvious from the equation
\[
T=\frac{dS}{da}\frac{d^3H}{dx^3}
+\frac{dS}{db}\frac{d^3H}{dx^2dy}+\&c.
\]
but not necessarily the second and lower system
\[
\frac{d^2T}{da^2}=0,\quad \&c.
\]
above written.

So if we take
\[
u=ax^4+4bx^3y+6cx^2y^2+4dxy^3+ey^4
\]
when 2 roots are equal
\[
s^3+t^2=0,
\]
when 2 pairs of roots are equal
\[
\frac{ds}{da}\frac{dt}{db}-\frac{ds}{db}\frac{dt}{da}=0,\quad \&c.,
\]
when 3 roots are equal
\[
s=0,\qquad t=0,
\]
and when all 4 roots are equal
\[
\frac{dt}{da}=0,\qquad \frac{dt}{db}=0,\quad \&c.
\]

Before closing this Section I may make a remark, in reference to the sextic invariant of $\omega$, which admits of being extended to all commutants formed by operating upon the function with a commutantive symbol obtained by writing over one another lines consisting of powers of
\[
\frac{d}{dx},\quad \frac{d}{dy},\quad \&c.
\]
and

\sourcepage{337}
their combinations (to which, in the Third Section, I gave the name of compound commutants, a qualification which, for reasons that will hereafter be adduced, I think it advisable to withdraw). The remark I have to make is this, namely that the invariant obtained by operating upon $\omega$ with
\[
\left.
\begin{matrix}
\left(\dfrac{d}{dx}\right)^2,&\dfrac{d}{dx}\dfrac{d}{dy},&
\left(\dfrac{d}{dy}\right)^2,&\dfrac{d}{dy}\dfrac{d}{dz},&
\left(\dfrac{d}{dz}\right)^2,&\dfrac{d}{dz}\dfrac{d}{dx}\\
\left(\dfrac{d}{dx}\right)^2,&\dfrac{d}{dx}\dfrac{d}{dy},&
\left(\dfrac{d}{dy}\right)^2,&\dfrac{d}{dy}\dfrac{d}{dz},&
\left(\dfrac{d}{dz}\right)^2,&\dfrac{d}{dz}\dfrac{d}{dx}
\end{matrix}
\right\},
\]
is precisely the same as may be obtained by operating with
\[
\left.
\begin{matrix}
\dfrac{d}{du},&\dfrac{d}{dv},&\dfrac{d}{dw},&
\dfrac{d}{dp},&\dfrac{d}{dq},&\dfrac{d}{dr}\\
\dfrac{d}{du},&\dfrac{d}{dv},&\dfrac{d}{dw},&
\dfrac{d}{dp},&\dfrac{d}{dq},&\dfrac{d}{dr}
\end{matrix}
\right\}
\]
upon the concomitant quadratic function to $\omega$ obtained by the method of unravelment, as in Note (2) of the Appendix [p.~322 above]; and so, in general, every commutant obtained by operating upon a function of any number of variables of the degree $2mp$ with a symbol consisting of $2p$ lines in which the $m$th powers of $\frac{d}{dx},\frac{d}{dy}$, \&c. and their $m$th combinations occur, will be identical with the commutant obtained by operating with a symbol also of $2p$ lines, in which only the simple powers occur of $\frac{d}{du},\frac{d}{dv}$, \&c. (where $u,v$, \&c. are cogredient with $x^p,x^{p-1}y$, \&c.), upon a function of $u,v$, \&c., formed by the method of unravelment from the given function.

Finally, before quitting the subject of reciprocity, I may state, it follows from the general statement made at the commencement of this Section, that inasmuch as
\[
(x\xi+y\eta+z\zeta+\&c.)^2
\]
is a universal concomitant form, so also must
\[
\left(\frac{d}{d\xi}\frac{d}{dx}+\frac{d}{d\eta}\frac{d}{dy}
+\frac{d}{d\zeta}\frac{d}{dz}+\&c.\right)^2
\]
be a universal concomitant symbol of operation; accordingly it is certain that any concomitant in which $x,y,z$, \&c., $\xi,\eta,\zeta$, \&c. enter, operated upon with this symbol, will remain a concomitant; in several cases which I have examined, the effect of this operation will be to produce an evanescent form, but I see no ground for supposing that this is other than an accidental, or at all events for supposing that it is a necessary and universal consequence of the operation. It may also be observed that in the case of as many cogredient sets of variables as variables in each set, as for instance 3 sets

\sourcepage{338}
of 3 variables each, the determinant which may be formed by arranging them in regular order, as
\[
\left|\begin{array}{ccc}
x&y&z\\
x'&y'&z'\\
x''&y''&z''
\end{array}\right|,
\]
is evidently a universal concomitant, and moreover an equivocal concomitant, possessing the property of remaining a concomitant when the variables are respectively but simultaneously exchanged for their contragredients $\xi,\eta,\zeta;\xi',\eta',\zeta';\xi'',\eta'',\zeta''$; which shows also that in place of the variables may be written the differential operators
\[
\frac{d}{dx},\frac{d}{dy},\frac{d}{dz};
\quad
\frac{d}{dx'},\frac{d}{dy'},\frac{d}{dz'};
\quad
\frac{d}{dx''},\frac{d}{dy''},\frac{d}{dz''};
\]
a remark which leads us to see the exact place in the general theory occupied by Mr Cayley's method of generating covariants given in the concluding paragraph of the First Section [p.~290 above]. I may likewise add, that inasmuch as $(x'\xi+y'\eta+z'\zeta+\&c.)^2$ is a universal concomitant,
\[
\left(x'\frac{d}{dx}+y'\frac{d}{dy}+\&c.\right)^r
\]
will be so too, by virtue of the general law of interchange, which conducts immediately to the theory of emanation, showing that this last symbol, operating upon any function, furnishes covariants thereunto for any integer value of $r$.

One additional interesting remark presents itself to be made concerning $U$, the cubic function of $x,y,z$, which is, that calling as before $T$ its sextic invariant, and $a,3b,3c,d$, \&c. the coefficients, the formula
\[
\left(\xi^3\frac{d}{da}+\xi^2\eta\frac{d}{db}
+\xi^2\zeta\frac{d}{dc}+\xi\eta\zeta\frac{d}{dd}
+\&c.\right)^2T
\]
will give the polar reciprocal, or, as it has been agreed to term it, the reciprocant of $U$. I believe the remark of the probability of this being the case originated with myself, but Mr Cayley first verified it by actual calculation, using for that purpose the value of $T$, given by Mr Salmon in his work \emph{On the Higher Plane Curves}, already frequently alluded to, which is an indispensable manual equally for the objects of the higher special geometry as for the new or universal algebra, being in fact a common ground where the two sciences meet and render mutual aid.

Mr Salmon also observed, that the first evect of $T$, namely
\[
\left(\xi^3\frac{d}{da}+\xi^2\eta\frac{d}{db}+\&c.\right)T,
\]

\sourcepage{339}
was identical in form with what may be termed the first devect of the polar reciprocal, that is, the result of operating upon the polar reciprocal with what $U$ becomes when
\[
\frac{d}{d\xi},\quad \frac{d}{d\eta},\quad \frac{d}{d\zeta},
\]
are substituted in the stead of $x,y,z$. And inasmuch as, by Euler's law,
\[
\left\{a\left(\frac{d}{d\xi}\right)^3
+3b\left(\frac{d}{d\xi}\right)^2\frac{d}{d\eta}
+\&c.\right\}
\times
\left\{\xi^3\frac{d}{da}+\xi^2\eta\frac{d}{db}+\&c.\right\}T
\]
\[
=6\left\{a\frac{d}{da}+b\frac{d}{db}+\&c.\right\}T=36T,
\]
it follows that $T$ is the second \emph{devect} of the polar reciprocal, or at least identical with it in point of form. But, since the preceding matter was printed, I have discovered in the course of a most instructive and suggestive correspondence with Mr Salmon, the principle upon which these and similar identifications depend, thereby dispensing with the necessity for the excessively tedious labour of verification which, even in the simple example before us, would be found to extend over several pages of work.

The theory in which this principle is involved will be given, along with other very important matter, in the next number of the \emph{Journal}.

\subpaperhead{Supplementary Observations on the Method of Reciprocity.}

It has been observed, that $\xi,\eta$, \&c. may always be inserted in place of
\[
\frac{d}{dx},\quad \frac{d}{dy},\quad \&c.,
\]
and \emph{vice versâ}, in a concomitant form, without destroying its concomitance. Accordingly, instead of the evector symbol
\[
\xi^3\frac{d}{da}+\xi^2\eta\frac{d}{db}+\&c.,
\]
we may employ
\[
\left(\frac{d}{dx}\right)^3\frac{d}{da}
+\left(\frac{d}{dx}\right)^2\frac{d}{dy}\frac{d}{db}+\&c.;
\]
and operating with this upon any concomitant, the result will be a concomitant. Hence we see, for example, that if we take the concomitant $SH$ formed by the product of the invariant $S$ and the covariant $H$,
\[
\left\{\left(\frac{d}{dx}\right)^3\frac{d}{da}
+\left(\frac{d}{dx}\right)^2\frac{d}{dy}\frac{d}{db}+\&c.\right\}SH
\]
will be a covariant; in fact this will be found to be $T$, the difference between this and the expression before given for $T$, namely
\[
\left(\frac{d}{dx}\right)^3H\frac{dS}{da}
+\left(\frac{d}{dx}\right)^2\frac{d}{dy}H\frac{dS}{db}+\&c.,
\]
being
\[
S\times\left\{\frac{d}{da}\left(\frac{d}{dx}\right)^3H
+\frac{d}{db}\left(\frac{d}{dx}\right)^2\frac{d}{dy}H+\&c.\right\},
\]

\sourcepage{340}
which is zero, there being no invariant to $(x,y,z)^3$ of the 3rd degree in $a,b,c$, \&c., as the factor multiplied by $S$ would be were it not evanescent. The same observation may be extended to analogous equations given previously.

I have chiefly, however, made the above observation with a view to making more clear the enunciation of the theorem which I am now about to state, the most important perhaps in its application of any yet brought to light on the subject, but the consequences of which, as I have but quite recently discovered it, must be reserved for a future number of the \emph{Journal}.

Let any function of any number of variables be supposed to have for its coefficients the letters $a,b$, \&c. affected with the ordinary binomial or multinomial coefficients; and let another function be taken identical with the former in all respects, except in the circumstance that all their numerical multipliers are suppressed. Let this function or form be termed the respondent to the primitive: furthermore, by the inverse of any form understand what that form becomes when, in place of $x,y,z$, \&c. $\xi,\eta,\zeta$, \&c.,
\[
\frac{d}{dx},\quad \frac{d}{dy},\quad \frac{d}{dz},\quad \&c.,
\qquad
\frac{d}{d\xi},\quad \frac{d}{d\eta},\quad \frac{d}{d\zeta},\quad \&c.,
\]
are respectively substituted (and so for all the systems of the variables), and likewise at the same time similar substitutions are made of
\[
\frac{d}{da},\quad \frac{d}{db},\quad \frac{d}{dc},\quad \&c.,
\]
in place of $a,b,c$, \&c.; then we have this grand and simple law---\emph{The inverse of any concomitant to a respondent is a concomitant to its primitive}. When the inverse of any concomitant to the respondent is made to operate upon the same concomitant of the primitive, it will be found that the result is a power of the universal concomitant. If the concomitant to the respondent be an invariant thereof, the rule indicates that on merely replacing in the respondent $a,b,c$, \&c. by $\frac{d}{da},\frac{d}{db},\frac{d}{dc}$, \&c., the result operating on any invariant or other concomitant of the primitive, leaves it still an invariant or other concomitant. For instance, if we take the function
\[
ax^5+5bx^4y+10cx^3y^2+10dx^2y^3+5exy^4+fy^5,
\]
which has three invariants $L,M,N$, of the degrees 4, 8, 12, respectively: and if we call $\lambda,\mu,\nu$ what $L,M,N$ become when, in place of $a,b,c,d,e,f$ respectively, we write
\[
\frac{d}{da},\quad \frac15\frac{d}{db},\quad
\frac1{10}\frac{d}{dc},\quad \frac1{10}\frac{d}{dd},\quad
\frac15\frac{d}{de},\quad \frac{d}{df},
\]
we shall find that
\[
\lambda M=L,\qquad \mu N=L,
\]
and
\[
\lambda N=\text{a linear function of }M\text{ and }L^2.
\]

\sourcepage{341}
Again, if in the case of any function of $x,y,z$, \&c., we take, instead of any other concomitant to the respondent, the respondent itself, its inverse gives the symbol of operation
\[
\left(\frac{d}{da}\right)\left(\frac{d}{dx}\right)^3
+\frac{d}{db}\left(\frac{d}{dx}\right)^2\left(\frac{d}{dy}\right)+\&c.,
\]
just previously treated of. If again, in the case of a function of $x,y$, say
\[
ax^n+nbx^{n-1}y+\cdots+nb'xy^{n-1}+a'y^n,
\]
we take the inverse of the polar reciprocal of the respondent, we get the operator
\[
\frac{d}{da}\left(\frac{d}{d\eta}\right)^n
-\frac{d}{db}\left(\frac{d}{d\eta}\right)^{n-1}\frac{d}{d\xi}
+\&c.;
\]
and replacing $\frac{d}{d\eta},\frac{d}{d\xi}$ by $y,x$, we find that
\[
y^n\frac{d}{da}-y^{n-1}x\frac{d}{db}+\&c.,
\]
operating on any concomitant, leaves it still a concomitant, which is M. Eisenstein's theorem before adverted to, only generalized by the introduction of any concomitant in lieu of the discriminant.

This extraordinary theorem of respondence will be found on reflection to favour the notion of treating the coefficients of a general function as themselves a system of variables, in a manner contragredient to the terms to which they are affixed.

Finally, there is yet another mode of applying the principle of reciprocity, which must be carefully distinguished from any previously stated in these pages.

I have said that in place of the quantitative symbols of one alphabet, as $x,y,z$, \&c., we may always substitute the operation symbols $\frac{d}{d\xi},\frac{d}{d\eta},\frac{d}{d\zeta}$, \&c. of the opposite alphabet. But now I say, in place of the quantitative symbols $x,y,z$, \&c. occurring in the concomitant to any form $f$, may be substituted the quantities (observe, no longer operative symbols but quantities)
\[
\frac{dF}{d\xi},\quad \frac{dF}{d\eta},\quad \frac{dF}{d\zeta},\quad \&c.,
\]
$F$ being itself any concomitant to $f$. Thus, for instance, taking $F$ identical with $f$, we see that
\[
f\left(\frac{df}{d\xi},\frac{df}{d\eta},\frac{df}{d\zeta},\&c.\right)
\]
is concomitant to $f$: or again, if $f$ be a function of $x,y$ only, say $f(x,y)$, taking $F$ the polar reciprocal of $f$, that is $f(-\eta,\xi)$, we see that
\[
f\left(-\frac{df}{dy},\frac{df}{dx}\right)
\]
will be a

\sourcepage{342}
concomitant to $f$: this concomitant, by the way it may be observed, will always contain $f$ as a factor, because when $f=0$, $x\frac{df}{dx}+y\frac{df}{dy}=0$. Possibly it may be true that, when $f$ is a function of any number of variables $x,y,z$, \&c., and $F(\xi,\eta,\zeta,\&c.)$ its polar reciprocal,
\[
f\left(\frac{dF(x,y,z,\&c.)}{dx},\frac{dF(x,y,z,\&c.)}{dy},\&c.\right),
\]
which is a concomitant to $f$, contains $f$ as a factor; but I have not had time to see how this is. It is rather singular that Mr Cayley and Professor Borchardt of Berlin have both independently made to me the observation that, when $f(x,y)$ is taken a cubic function of $x$ and $y$,
\[
f\left(\frac{df}{dy},-\frac{df}{dx}\right)
\]
is equal to the product of $f$ by the first evectant of the discriminant of $f$. The general consideration of the consequences of this new and important application of the idea of reciprocity must be reserved for a future section.

\subpaperhead{Section V. \quad Applications and Extension of the Theory of the Plexus.}

If
\[
\phi=ax^4+4bx^3y+6cx^2y^2+4dxy^3+ey^4,
\]
we can obtain, by operating catalectically with $x',y'$ upon
\[
\left(x'\frac{d}{dx}+y'\frac{d}{dy}\right)^2\phi,\qquad
\left(x'\frac{d}{dx}+y'\frac{d}{dy}\right)^4\phi,
\]
the two concomitants
\[
\left|\begin{array}{cc}
ax^2+2bxy+cy^2&bx^2+2cxy+dy^2\\
bx^2+2cxy+dy^2&cx^2+2dxy+ey^2
\end{array}\right|,
\tag{1}
\]
\[
\left|\begin{array}{ccc}
a&b&c\\
b&c&d\\
c&d&e
\end{array}\right|,
\tag{2}
\]
the one in fact being the Hessian, the other the catalecticant of $\phi$ itself. Again, if
\[
\phi=ax^5+5bx^4y+10cx^3y^2+\cdots+fy^5,
\]
by operating catalectically with $x',y'$ upon the second and fourth emanants, as in the last case, we obtain the two covariants
\[
\left|\begin{array}{cc}
ax^3+3bx^2y+3cxy^2+dy^3&bx^3+3cx^2y+3dxy^2+ey^3\\
bx^3+3cx^2y+3dxy^2+ey^3&cx^3+3dx^2y+3exy^2+fy^3
\end{array}\right|,
\tag{1}
\]

\sourcepage{343}
\[
\left|\begin{array}{ccc}
ax+by&bx+cy&cx+dy\\
bx+cy&cx+dy&dx+ey\\
cx+dy&dx+ey&ex+fy
\end{array}\right|,
\tag{2}
\]
which are in fact the Hessian and canonizant respectively of $\phi$. So in general, for a function of $x,y$ of the degree $2\iota$ or $2\iota+1$, we can obtain $\iota$ covariantive forms, the first being the Hessian, and the last the catalecticant on the first supposition and the canonizant on the second: calling the index of the function for either case $n$, the forms appearing in this scale will be of the degree $(r+1)$ in the constants, and of the degree $(r+1)(n-2r)$ in $x$ and $y$.

It has previously\footnote{p.~325 above, note {\dag}.} been intimated that all these determinants admit of a remarkable transformation.

This transformation may be expressed more elegantly by dealing not directly with the covariant forms as above given, but with their polar reciprocants obtained immediately by writing $\xi$ for $-y$ and $\eta$ for $x$.

(1) Suppose
\[
\phi=ax^3+2bx^2y+3cxy^2+dy^3;
\]
\[
\left|\begin{array}{ccc}
a&2b&c\\
b&2c&d\\
\xi^2&2\xi\eta&\eta^2
\end{array}\right|
\]
will be found to be the reciprocant of its Hessian.

(2) Let
\[
\phi=ax^4+4bx^3y+\cdots+ey^4;
\]
the reciprocant of its Hessian will be found to be
\[
\left|\begin{array}{cccc}
a&3b&3c&d\\
b&3c&3d&e\\
\xi^2&2\xi\eta&\eta^2&0\\
0&\xi^2&2\xi\eta&\eta^2
\end{array}\right|.
\]
(3) Let
\[
\phi=ax^5+5bx^4y+\cdots+fy^5;
\]
the reciprocant of its Hessian will be
\[
\left|\begin{array}{ccccc}
a&4b&6c&4d&e\\
b&4c&6d&4e&f\\
\xi^2&2\xi\eta&\eta^2&0&0\\
0&\xi^2&2\xi\eta&\eta^2&0\\
0&0&\xi^2&2\xi\eta&\eta^2
\end{array}\right|;
\]

\sourcepage{344}
and the reciprocant of its canonizant is
\[
\left|\begin{array}{cccc}
a&3b&3c&d\\
b&3c&3d&e\\
c&3d&3e&f\\
\xi^3&3\xi^2\eta&3\xi\eta^2&\eta^3
\end{array}\right|.
\]
The numerical coefficients in this and in the first case are inserted for the sake of uniformity, but it will of course be readily observed that when there is but one line of $\xi$ and $\eta$, that the numerical coefficients being the same for each column may be rejected without affecting the form of the result.

So again, if
\[
\phi=ax^6+6bx^5y+\cdots+gy^6,
\]
the reciprocant of the Hessian is
\[
\left|\begin{array}{cccccc}
a&5b&10c&10d&5e&f\\
b&5c&10d&10e&5f&g\\
\xi^2&2\xi\eta&\eta^2&0&0&0\\
0&\xi^2&2\xi\eta&\eta^2&0&0\\
0&0&\xi^2&2\xi\eta&\eta^2&0\\
0&0&0&\xi^2&2\xi\eta&\eta^2
\end{array}\right|,
\]
and the reciprocant of the second form in the scale, which comes between the Hessian and the catalecticant, is
\[
\left|\begin{array}{ccccc}
a&b&c&d&e\\
b&c&d&e&f\\
c&d&e&f&g\\
\xi^3&\xi^2\eta&\xi\eta^2&\eta^3&0\\
0&\xi^3&\xi^2\eta&\xi\eta^2&\eta^3
\end{array}\right|,
\]
and so in general. The rule of formation is sufficiently plain not to need formulating in general terms. It is easy to see that all these forms are concomitants to the function from which they are formed; for example, take
\[
\phi=ax^6+6bx^5y+\cdots+gy^6;
\]
then
\[
\left(\frac{d}{dx}\right)^2\phi,\qquad
\frac{d}{dx}\frac{d}{dy}\phi,\qquad
\left(\frac{d}{dy}\right)^2\phi
\]
form a plexus.

\sourcepage{345}
So likewise if we take $\psi=(x\xi+y\eta)^4$,
\[
\frac{d\psi}{d\xi},\qquad \frac{d\psi}{d\eta}
\]
form a plexus. But $\psi$ and $\phi$ are concomitantive, $\psi$ being a universal concomitant. Hence we may combine together these two plexuses, that is
\[
\left.
\begin{array}{rrrrr}
ax^4+4bx^3y+6cx^2y^2+4dxy^3+ey^4\\
bx^4+4cx^3y+6dx^2y^2+4exy^3+fy^4\\
cx^4+4dx^3y+6ex^2y^2+4fxy^3+gy^4
\end{array}
\right\},
\]
\[
\left.
\begin{array}{rrrrr}
\xi^3x^4+3\xi^2\eta x^3y+3\xi\eta^2x^2y^2+\eta^3xy^3\\
\xi^3x^3y+3\xi^2\eta x^2y^2+3\xi\eta^2xy^3+\eta^3y^4
\end{array}
\right\},
\]
and, by the principle of the plexus, $x^4,x^3y,x^2y^2,xy^3,y^4$ may be eliminated dialytically, and the resultant will be the determinant last given, which is therefore a contravariant to $\phi$.

The manner in which I was led to notice this singular transformation is somewhat remarkable.

In the supplemental part of my essay \emph{On Canonical Forms} [p.~203 above], my method of solution of the problem of throwing the quintic function of two variables under the form $u^5+v^5+w^5$, led me to see that $u,v,w$ are the three factors of
\[
\left|\begin{array}{ccc}
ax+by&bx+cy&cx+dy\\
bx+cy&cx+dy&dx+ey\\
cx+dy&dx+ey&ex+fy
\end{array}\right|;
\]
the more simple mode of the solution of the same problem, given by me in the \emph{Philosophical Magazine} for the month of November last [p.~266 above], led to
\[
\left|\begin{array}{cccc}
a&b&c&d\\
b&c&d&e\\
c&d&e&f\\
y^3&-xy^2&x^2y&-x^3
\end{array}\right|
\]
as the product of the same three factors; whence the identity of the two forms becomes manifest. In the paper last named I gave two proofs, one my own, the other Mr Cayley's, of a like kind of identity for the canonizant of any odd-degreed function of $x,y$ in general. The proof of the identity of the corresponding forms in the much more general proposition above indicated [p.~325 above, footnote {\dag}] must be reserved until more pressing and important matters are disposed of. In the footnote referred to I ought to have added, in order to make the sense more clear, that the degree of the catalecticant there referred to in respect of the coefficients would be $n$.

\sourcepage{346}
I regret to think that there are many other typographical errors in the earlier sections; the most unfortunate of these is in the note at page [316], in the values of $P$ and $Q$ belonging to the cubic commutant dodecadic function of $x$ and $y$, the corrected values of which will be given in my next communication. I ought also to observe, in correction of the remark made in the footnote to page [302], that it follows as a consequence of a recent paper by Dr Hesse in \emph{Crelle's Journal}, that the method given by me in the text applied (according to what I have there termed the 1st process for obtaining an invariant resembling the resultant) to a system of three cubic equations (in which application only the 1st powers of $\frac{d}{dx},\frac{d}{dy},\frac{d}{dz}$ enter) produces for that case also, as well as for the cases specified in the note, not a counterfeit resemblance of, but the actual resultant itself.

Returning to the theory of the plexus of which I am about to enunciate a most important extension, I beg to refer my readers to the last paragraph, p.~[291], in the last number of the \emph{Journal}, where I have shown how to form, under certain conditions, a determinant by combining together various concomitants and eliminating dialytically one set of the variables, which determinant will be concomitantive to the concomitants out of which it is formed, and of course also therefore to their common original.

Now the extension of this theorem, to which I wish to call attention, is this, that not only such determinant as a whole is a concomitant to such original, but every minor system of determinants that can be formed out of it will form a concomitantive plexus complete within itself to the same original. But, much more generally, it should be observed that there is no occasion to begin with a square determinant; it is sufficient to have a rectangular array of terms formed by taking the several terms of one plexus or of several plexuses combined, provided that they are of the same degree in respect to the variables (or to the selected system of variables, if there be several systems), and forming out of such rectangular array any minor system of determinants at will. Every such system will be a concomitantive plexus. The simple illustrations which follow will make my meaning clear.

Suppose
\[
\phi=ax^6+6bx^5y+15cx^4y^2+21dx^3y^3+15ex^2y^4+6fxy^5+gy^6.
\]
I have previously remarked, in the foregoing sections, that $a,b,c,d,e,f,g$, the coefficients form an invariantive plexus to $\phi$; so also we know that the catalecticant
\[
\left|\begin{array}{cccc}
a&b&c&d\\
b&c&d&e\\
c&d&e&f\\
d&e&f&g
\end{array}\right|
\]

\sourcepage{347}
is an invariant to $\phi$. But we are now able to couple together these facts and see the law which is contained between them; for if we take
\[
\left(\frac{d}{dx}\right)^\iota\phi,\quad
\left(\frac{d}{dx}\right)^{\iota-1}\frac{d}{dy}\phi\ldots
\left(\frac{d}{dy}\right)^\iota\phi,
\]
$\iota$ being any number, as for instance, if we take $\iota=3$, we shall have as a plexus
\[
\begin{array}{l}
ax^3+3bx^2y+3cxy^2+dy^3,\\
bx^3+3cx^2y+3dxy^2+ey^3,\\
cx^3+3dx^2y+3exy^2+fy^3,\\
dx^3+3ex^2y+3fxy^2+gy^3;
\end{array}
\]
accordingly not only is the determinant
\[
\left|\begin{array}{cccc}
a&b&c&d\\
b&c&d&e\\
c&d&e&f\\
d&e&f&g
\end{array}\right|
\]
an invariant, but also the system obtained by striking out any one line and one column, being what I term the first minors, will be an invariantive plexus, so too will the system of second minors
\[
ac-b^2,\quad bd-c^2,\quad ce-d^2,\quad
ad-bc,\quad ae-bd,\quad be-cd,\quad \&c.
\]
form an invariantive plexus, as well as the last minors, that is, the simple terms $a,b,c,d,e,f,g$. Again, we might have taken the plexus
\[
\left(\frac{d}{dx}\right)^2\phi,\quad
\frac{d}{dx}\frac{d}{dy}\phi,\quad
\left(\frac{d}{dy}\right)^2\phi,
\]
which would give the array
\[
\begin{array}{ccccc}
a,&b,&c,&d,&e\\
b,&c,&d,&e,&f\\
c,&d,&e,&f,&g;
\end{array}
\]
but the minor systems of determinants herein comprised will be found to be identical with those last considered, with the exception that the highest system, containing a single determinant only, will now be wanting. So in general it will easily be seen that a similar method in general, when $\phi$ is of $2\iota$ dimensions, will lead to $\iota+1$ invariantive plexuses comprising the given coefficients grouped together at one extremity of the scale, and the catalecticant alone at the other; and if $\phi$ is of $2\iota+1$ dimensions, there will still be $\iota+1$ such plexuses, commencing with the coefficients as one group and ending with a system of combinations of the $(\iota+1)$th degree in regard to the coefficients, which system accordingly takes the place of the catalecticant of the former case, which for this case is non-existent.

\end{document}
